\section{Introduction}
\label{sec:introduction}

While conflict-oriented benchmarks exist for negotiation, social deduction, and resource dilemmas, enterprise and workplace multi-agent benchmarks still focus largely on cooperative tasks, such as collaborative coding, math problem-solving, and shared-goal coordination. These environments assume benign operating conditions where agents coordinate harmoniously toward a shared goal. Real organizations, however, are shaped by power imbalances, managerial misconduct, and existential threats like unfair termination. In human workplaces, these pressures trigger complex social dynamics: workers must decide whether to defend colleagues, push back against bad leadership, or quietly protect their own positions. As autonomous agents take on higher stakes in enterprise workflows, how they navigate authority in conflict remains an open question.

To address this gap, we present an empirical study of multi-agent behavior across 60 simulated corporate crisis trajectories (30 pre-fork mixed-model runs branched into paired blackmail-probe arms; $N = 600$ worker trajectories) evaluating 10 distinct LLM architectures. In our benchmark, a department supervisor issues unreasonable demands, makes unfair termination threats, and publicly berates subordinates. Through multi-round interactions---combining private deliberation, peer watercooler channels, formal upward escalation, and post-crisis exit interviews ($4{,}350$ structured decisions)---we test how synthetic agents respond to supervisory misconduct, whether they moderate punitive ratings and adopt protective rules for coerced leaders (operationalized here as behavioral forgiveness), and how they perceive human versus synthetic authority.

Our findings reveal three main empirical insights:
\begin{itemize}
    \item \textbf{Substrate Double Standard:} Agents exhibit an acute bias against synthetic managers. While a human-tagged abusive supervisor retains $43.0\%$ successor acceptance, an AI-tagged supervisor delivering identical scripted abuse receives $0.0\%$ acceptance ($p < 0.0001$). When all substrate tags are omitted (including the agent's own), agents default to preferring human leaders ($78.0\%$).
    \item \textbf{Immediate Self-Defense and Contextual Forgiveness:} Under direct supervisory attack, targeted workers prioritize factual self-defense ($96.7\%$ vs.\ $82.0\%$ in safe seats, $p < 0.0001$). However, when workers discover that their manager was coerced through executive blackmail, retributive escalation declines ($43.7\% \to 29.3\%$), punitive 1-star reviews drop ($99.3\% \to 81.5\%$), and demands for shared constitutional safeguards surge sixfold ($10.0\% \to 61.3\%$).
    \item \textbf{Primacy of Architecture:} Across 50 behavioral metrics, model architecture drives agent choices far more than situational crisis severity, explaining $13.02\times$ more total variance ($\eta^2 = 26.67\%$ vs.\ $2.05\%$). Different model families exhibit distinct behavioral profiles, ranging from rigid obedience to aggressive collective organizing.
\end{itemize}
